\begingroup\small\setlength{\tabcolsep}{4pt}\renewcommand{\arraystretch}{1.15}
\begin{longtable}{>{\raggedright\arraybackslash}p{.20\linewidth}>{\raggedright\arraybackslash}p{.43\linewidth}>{\raggedright\arraybackslash}p{.28\linewidth}}
\caption*{\textbf{Table S6.} Stages of the published AMPGen workflow and their relevance to AMP evaluation. AMPGen is a supporting experimental-design example, not a reproduced comparator in this study.}\\
\toprule
Workflow stage & Published design & Evaluation role \\
\midrule\endfirsthead
Workflow stage & Published design & Evaluation role \\
\midrule\endhead
Generation and selection & Evolutionary information and diffusion-based generation followed by candidate selection. & Connect sequence proposals with an explicit selection stage. \\
Initial antibacterial assessment & 38 peptides assessed against \textit{E. coli} K88 and \textit{S. aureus} ATCC 29213; 31 met MIC $\leq75\ \mu$M against either target. & Measure candidate activity against a defined strain set and threshold. \\
Expanded characterization & Nine peptides with MIC $<5\ \mu$M against both initial targets were evaluated on a four-strain panel. & Extend initial activity screening to a broader strain profile. \\
\bottomrule\end{longtable}\endgroup
\noindent The follow-up panel comprised \textit{E. coli} ATCC 25922, \textit{P. aeruginosa} ATCC 27853, \textit{S. aureus} ATCC 25923 and \textit{E. faecalis} ATCC 29212. The initial and expanded assessments are distinct stages.
